\documentclass[aspectratio=169]{beamer}
\usetheme{metropolis}
\usepackage{graphicx}
\usepackage{tikz}
\usepackage{amsmath}
\usepackage{listings}
\usepackage{xcolor}
\usepackage{booktabs}
\usepackage{emoji}

% Code listing setup
\lstset{
  basicstyle=\ttfamily\small,
  keywordstyle=\color{blue},
  commentstyle=\color{gray},
  stringstyle=\color{red},
  showstringspaces=false,
  breaklines=true,
  frame=single,
  language=Python
}

\title{Word Embeddings: Word2Vec, GloVe, FastText}
\subtitle{Lecture 8: The Neural Revolution in NLP}
\author{LLM Course - Week 3}
\date{Winter 2026}

\begin{document}

% ============================================================
% TITLE SLIDE
% ============================================================
\begin{frame}
\titlepage
\end{frame}

% ============================================================
% OVERVIEW
% ============================================================
\begin{frame}{Today's Lecture 📋}
\large
\begin{enumerate}
    \item 🚀 \textbf{The 2013 Revolution: Word2Vec}
    \item 🏗️ \textbf{Word2Vec Architectures: CBOW \& Skip-gram}
    \item 🌍 \textbf{GloVe: Global Vectors}
    \item 🔤 \textbf{FastText: Subword Information}
    \item 📊 \textbf{Comparison \& Evaluation}
    \item 💡 \textbf{Applications \& Best Practices}
\end{enumerate}

\vspace{1em}
\textit{Goal: Understand how neural methods learn dense word representations}
\end{frame}

% ============================================================
% FROM COUNT TO PREDICT
% ============================================================
\begin{frame}{From Count-Based to Prediction-Based 🎯}
\begin{columns}[T]
\begin{column}{0.48\textwidth}
\textbf{Count-Based (LSA, LDA):}
\begin{itemize}
    \item Build co-occurrence matrix
    \item Apply matrix factorization
    \item Global statistics
    \item Linear relationships
    \item Interpretable factors
\end{itemize}

\vspace{0.5em}
\textbf{Advantages:}
\begin{itemize}
    \item Fast training
    \item Leverages global statistics
    \item Mathematically grounded
\end{itemize}
\end{column}

\pause

\begin{column}{0.48\textwidth}
\textbf{Prediction-Based (Word2Vec):}
\begin{itemize}
    \item Predict context from word
    \item Train neural network
    \item Local context windows
    \item Non-linear relationships
    \item Learned representations
\end{itemize}

\vspace{0.5em}
\textbf{Advantages:}
\begin{itemize}
    \item Better semantic capture
    \item Scalable to huge corpora
    \item Dense, low-dimensional
\end{itemize}
\end{column}
\end{columns}

\vspace{1em}
\begin{center}
\alert{2013: Word2Vec showed prediction-based methods could outperform count-based!}
\end{center}
\end{frame}

% ============================================================
% WORD2VEC INTRO
% ============================================================
\begin{frame}{Word2Vec: The Revolution 🚀}
\begin{center}
\LARGE
\textbf{2013: Everything changed}
\end{center}

\vspace{1em}

\begin{columns}[T]
\begin{column}{0.48\textwidth}
\textbf{Key Innovation:}
\begin{itemize}
    \item Shallow neural network
    \item Predict context from word (CBOW)
    \item Predict word from context (Skip-gram)
    \item Dense, low-dimensional vectors
    \item Captures semantic relationships
    \item \alert{Fast} training with negative sampling
\end{itemize}

\vspace{0.5em}
\textbf{Impact:}
\begin{itemize}
    \item Sparked deep learning in NLP
    \item Pre-training became standard
    \item Enabled transfer learning
\end{itemize}
\end{column}

\begin{column}{0.48\textwidth}
\textbf{Famous Results:}
\begin{center}
\Large
$\vec{king} - \vec{man} + \vec{woman} \approx \vec{queen}$
\end{center}

\vspace{1em}
\textbf{Other Examples:}
\begin{itemize}
    \item $\vec{Paris} - \vec{France} + \vec{Italy} \approx \vec{Rome}$
    \item $\vec{walking} - \vec{walk} + \vec{swim} \approx \vec{swimming}$
    \item $\vec{bigger} - \vec{big} + \vec{small} \approx \vec{smaller}$
\end{itemize}

\vspace{0.5em}
\alert{Vector arithmetic captures semantic relationships!}
\end{column}
\end{columns}

\vspace{0.5em}
\small
\textit{Reference: Mikolov et al. (2013a). "Efficient Estimation of Word Representations in Vector Space"}
\end{frame}

% ============================================================
% WORD2VEC INTUITION
% ============================================================
\begin{frame}{Word2Vec: Core Intuition 💡}
\begin{center}
\Large
\textit{Words that appear in similar contexts should have similar representations}
\end{center}

\vspace{1em}

\begin{exampleblock}{Example Context Window}
\begin{center}
"The quick brown \textbf{[TARGET]} jumped over the lazy dog"
\end{center}

\textbf{Context:} [The, quick, brown, jumped, over, the, lazy, dog]

\textbf{Target:} fox
\end{exampleblock}

\vspace{1em}

\textbf{Key Idea:}
\begin{itemize}
    \item Train a model to predict target from context (or vice versa)
    \item The learned \alert{hidden layer weights} become our word vectors!
    \item Similar words will have similar weight patterns
\end{itemize}

\vspace{0.5em}
\begin{center}
\textit{We don't actually care about the prediction task—we want the embeddings!}
\end{center}
\end{frame}

% ============================================================
% CBOW
% ============================================================
\begin{frame}{CBOW: Continuous Bag of Words 🎒}
\textbf{Predict the center word from context words}

\vspace{1em}

\begin{columns}[T]
\begin{column}{0.48\textwidth}
\begin{center}
\begin{tikzpicture}[scale=0.8]
    % Input layer
    \node at (-2, 4.5) {the};
    \node at (-1, 4.5) {cat};
    \node at (1, 4.5) {on};
    \node at (2, 4.5) {the};

    % Input vectors
    \foreach \x in {-2,-1,1,2}
        \draw[fill=blue!20] (\x-0.2, 3.7) rectangle (\x+0.2, 4.2);

    % Arrows to projection layer
    \foreach \x in {-2,-1,1,2}
        \draw[->] (\x, 3.7) -- (0, 2.5);

    % Projection/Hidden layer
    \draw[fill=green!20] (-0.8,1.8) rectangle (0.8,2.5);
    \node at (0, 2.15) {Hidden};
    \node[font=\tiny] at (0, 1.5) {Average/Sum};

    % Arrow to output
    \draw[->] (0, 1.8) -- (0, 0.8);

    % Output layer
    \draw[fill=red!20] (-0.5,0.3) rectangle (0.5,0.8);
    \node at (0, -0.1) {\textbf{sat}};

    % Labels
    \node[font=\small] at (-3.5, 4.5) {Input:};
    \node[font=\small] at (-3.5, 2.15) {Hidden:};
    \node[font=\small] at (-3.5, 0.5) {Output:};
\end{tikzpicture}
\end{center}
\end{column}

\begin{column}{0.48\textwidth}
\textbf{Architecture:}
\begin{enumerate}
    \item Input: One-hot vectors of context words
    \item Hidden: Average of input embeddings
    \item Output: Softmax over vocabulary
\end{enumerate}

\vspace{0.5em}
\textbf{Training:}
\begin{itemize}
    \item Context window size (e.g., 5 words)
    \item Maximize $P(\text{center}|\text{context})$
    \item Backpropagation updates embeddings
\end{itemize}

\vspace{0.5em}
\textbf{Characteristics:}
\begin{itemize}
    \item \alert{Faster} to train than Skip-gram
    \item Better for frequent words
    \item Smooths over context
    \item Good for syntactic tasks
\end{itemize}
\end{column}
\end{columns}

\vspace{0.5em}
\textbf{Objective:} $\max \frac{1}{T}\sum_{t=1}^T \log P(w_t | w_{t-c},...,w_{t+c})$
\end{frame}

% ============================================================
% SKIP-GRAM
% ============================================================
\begin{frame}{Skip-gram: The Inverse Approach 🔄}
\textbf{Predict context words from the center word}

\vspace{1em}

\begin{columns}[T]
\begin{column}{0.48\textwidth}
\begin{center}
\begin{tikzpicture}[scale=0.8]
    % Input layer
    \node at (0, 4.5) {\textbf{sat}};

    % Input vector
    \draw[fill=blue!20] (-0.3, 3.7) rectangle (0.3, 4.2);

    % Arrow down
    \draw[->] (0, 3.7) -- (0, 2.5);

    % Hidden layer
    \draw[fill=green!20] (-0.8,1.8) rectangle (0.8,2.5);
    \node at (0, 2.15) {Hidden};

    % Arrows to output
    \foreach \x in {-2,-1,1,2}
        \draw[->] (0, 1.8) -- (\x, 0.8);

    % Output layer
    \foreach \x in {-2,-1,1,2}
        \draw[fill=red!20] (\x-0.2, 0.3) rectangle (\x+0.2, 0.8);

    \node at (-2, 0) {the};
    \node at (-1, 0) {cat};
    \node at (1, 0) {on};
    \node at (2, 0) {the};

    % Labels
    \node[font=\small] at (-3.5, 4.5) {Input:};
    \node[font=\small] at (-3.5, 2.15) {Hidden:};
    \node[font=\small] at (-3.5, 0.5) {Output:};
\end{tikzpicture}
\end{center}
\end{column}

\begin{column}{0.48\textwidth}
\textbf{Architecture:}
\begin{enumerate}
    \item Input: One-hot vector of center word
    \item Hidden: Word embedding (lookup)
    \item Output: Softmax over vocabulary (×C times)
\end{enumerate}

\vspace{0.5em}
\textbf{Training:}
\begin{itemize}
    \item For each context position
    \item Maximize $P(\text{context}|\text{center})$
    \item Independent predictions
\end{itemize}

\vspace{0.5em}
\textbf{Characteristics:}
\begin{itemize}
    \item \alert{Slower} than CBOW
    \item Better for rare words
    \item More training examples per word
    \item Better semantic representations
\end{itemize}
\end{column}
\end{columns}

\vspace{0.5em}
\textbf{Objective:} $\max \frac{1}{T}\sum_{t=1}^T \sum_{-c \leq j \leq c, j \neq 0} \log P(w_{t+j} | w_t)$
\end{frame}

% ============================================================
% NEGATIVE SAMPLING
% ============================================================
\begin{frame}{Negative Sampling: The Speed Trick ⚡}
\textbf{Problem:} Softmax over entire vocabulary is too expensive!

\begin{align*}
P(w_O|w_I) = \frac{\exp(v_{w_O}^T v_{w_I})}{\sum_{w=1}^V \exp(v_w^T v_{w_I})}
\end{align*}

\vspace{0.5em}
For 100k vocabulary: Need to compute 100k exponentials per training example! 😱

\pause
\vspace{1em}

\textbf{Solution: Negative Sampling}

Instead of predicting across all words, create a binary classification:
\begin{itemize}
    \item \textbf{Positive sample:} Actual context word (label = 1)
    \item \textbf{Negative samples:} K random words (label = 0)
\end{itemize}

\vspace{0.5em}

\begin{align*}
\log \sigma(v_{w_O}^T v_{w_I}) + \sum_{i=1}^k \mathbb{E}_{w_i \sim P_n(w)} [\log \sigma(-v_{w_i}^T v_{w_I})]
\end{align*}

\vspace{0.5em}
\textbf{Typical K:} 5-20 for large datasets, 2-5 for small datasets

\vspace{0.5em}
\alert{Result:} Training becomes $O(k)$ instead of $O(V)$ per example!

\small
\textit{Reference: Mikolov et al. (2013b). "Distributed Representations of Words and Phrases"}
\end{frame}

% ============================================================
% WORD2VEC CODE
% ============================================================
\begin{frame}[fragile]{Word2Vec in Practice 💻}
\begin{lstlisting}[language=Python]
from gensim.models import Word2Vec
import gensim.downloader as api

# Load pre-trained model (300 dimensions, 3B words)
model = api.load('word2vec-google-news-300')

# Find similar words
similar = model.most_similar('computer', topn=5)
print(similar)
# Output: [('computers', 0.72), ('laptop', 0.69), ('PC', 0.68), ...]

# Word analogies: king - man + woman = ?
result = model.most_similar(
    positive=['king', 'woman'],
    negative=['man'],
    topn=1
)
print(result)
# Output: [('queen', 0.71)]

# Train your own model
sentences = [
    ['the', 'cat', 'sat', 'on', 'the', 'mat'],
    ['the', 'dog', 'ran', 'in', 'the', 'park'],
    # ... more sentences
]

custom_model = Word2Vec(
    sentences,
    vector_size=100,    # embedding dimension
    window=5,           # context window
    min_count=1,        # ignore rare words
    sg=1,               # 1=skip-gram, 0=CBOW
    negative=5,         # negative sampling
    workers=4           # parallel threads
)
\end{lstlisting}
\end{frame}

% ============================================================
% GLOVE INTRO
% ============================================================
\begin{frame}{GloVe: Global Vectors for Word Representation 🌍}
\textbf{Combining the best of count-based and prediction-based methods}

\vspace{1em}

\begin{columns}[T]
\begin{column}{0.48\textwidth}
\textbf{Motivation:}
\begin{itemize}
    \item Word2Vec uses \alert{local} context windows
    \item LSA uses \alert{global} co-occurrence statistics
    \item Can we get the best of both?
\end{itemize}

\vspace{0.5em}
\textbf{Key Insight:}

Ratios of co-occurrence probabilities encode meaning better than raw probabilities!

\vspace{0.5em}
\begin{center}
\small
\begin{tabular}{lcc}
\toprule
$P(k|w)$ & ice & steam \\
\midrule
solid & high & low \\
gas & low & high \\
water & med & med \\
fashion & low & low \\
\bottomrule
\end{tabular}
\end{center}
\end{column}

\begin{column}{0.48\textwidth}
\textbf{Ratio Encodes Relevance:}

\vspace{0.5em}
\begin{align*}
\frac{P(\text{solid}|\text{ice})}{P(\text{solid}|\text{steam})} &\gg 1 \\
\frac{P(\text{gas}|\text{ice})}{P(\text{gas}|\text{steam})} &\ll 1 \\
\frac{P(\text{water}|\text{ice})}{P(\text{water}|\text{steam})} &\approx 1
\end{align*}

\vspace{0.5em}
\textbf{Idea:}
Learn word vectors such that their dot product relates to co-occurrence probability ratios

\vspace{0.5em}
$\vec{w}_i^T \vec{w}_j \approx \log P(i,j)$
\end{column}
\end{columns}

\vspace{0.5em}
\small
\textit{Reference: Pennington et al. (2014). "GloVe: Global Vectors for Word Representation"}
\end{frame}

% ============================================================
% GLOVE OBJECTIVE
% ============================================================
\begin{frame}{GloVe: The Objective Function 🎯}
\textbf{Goal:} Learn vectors that capture co-occurrence statistics

\vspace{1em}

\textbf{Objective Function:}
\begin{align*}
J = \sum_{i,j=1}^{V} f(X_{ij}) \left(\vec{w}_i^T \tilde{\vec{w}}_j + b_i + \tilde{b}_j - \log X_{ij}\right)^2
\end{align*}

where:
\begin{itemize}
    \item $X_{ij}$ = number of times word $j$ appears in context of word $i$
    \item $\vec{w}_i$ = word vector for word $i$
    \item $\tilde{\vec{w}}_j$ = separate context vector for word $j$
    \item $b_i, \tilde{b}_j$ = bias terms
    \item $f(X_{ij})$ = weighting function
\end{itemize}

\vspace{1em}

\textbf{Weighting Function:}
\begin{align*}
f(x) = \begin{cases}
(x/x_{\max})^\alpha & \text{if } x < x_{\max} \\
1 & \text{otherwise}
\end{cases}
\end{align*}

\textbf{Purpose:} Prevent very common co-occurrences from dominating (typically $x_{\max}=100$, $\alpha=0.75$)
\end{frame}

% ============================================================
% GLOVE VS WORD2VEC
% ============================================================
\begin{frame}{GloVe vs. Word2Vec 🥊}
\begin{columns}[T]
\begin{column}{0.48\textwidth}
\textbf{Word2Vec (Skip-gram):}
\begin{itemize}
    \item Local context windows
    \item Online training
    \item Predicts context from word
    \item Stochastic updates
    \item Negative sampling trick
    \item Each occurrence matters
\end{itemize}

\vspace{0.5em}
\textbf{Advantages:}
\begin{itemize}
    \item Works well on small corpora
    \item Can train incrementally
    \item Captures local patterns
\end{itemize}
\end{column}

\begin{column}{0.48\textwidth}
\textbf{GloVe:}
\begin{itemize}
    \item Global co-occurrence matrix
    \item Batch training
    \item Factorizes co-occurrence matrix
    \item Deterministic
    \item Weighted least squares
    \item Aggregates all occurrences
\end{itemize}

\vspace{0.5em}
\textbf{Advantages:}
\begin{itemize}
    \item Faster training (large corpora)
    \item Better statistical efficiency
    \item More interpretable
\end{itemize}
\end{column}
\end{columns}

\vspace{1em}

\begin{block}{Performance}
In practice: Similar performance on most tasks! Choose based on:
\begin{itemize}
    \item Corpus size (GloVe better for large)
    \item Training infrastructure (GloVe needs memory for matrix)
    \item Availability of pre-trained models
\end{itemize}
\end{block}
\end{frame}

% ============================================================
% GLOVE CODE
% ============================================================
\begin{frame}[fragile]{GloVe in Practice 💻}
\begin{lstlisting}[language=Python]
import gensim.downloader as api
import numpy as np

# Load pre-trained GloVe embeddings
# Options: glove-wiki-gigaword-50, -100, -200, -300
# Or: glove-twitter-25, -50, -100, -200
glove = api.load("glove-wiki-gigaword-100")

# Use just like Word2Vec
similar = glove.most_similar('computer', topn=5)
print(similar)

# Analogies
result = glove.most_similar(
    positive=['france', 'berlin'],
    negative=['paris'],
    topn=1
)
print(result)  # Should be close to 'germany'

# Vector arithmetic
king = glove['king']
man = glove['man']
woman = glove['woman']
result_vec = king - man + woman

# Find closest word
closest = glove.similar_by_vector(result_vec, topn=1)
print(closest)  # Should be close to 'queen'

# Compute similarity
similarity = glove.similarity('cat', 'dog')
print(f"Similarity: {similarity:.3f}")
\end{lstlisting}
\end{frame}

% ============================================================
% FASTTEXT INTRO
% ============================================================
\begin{frame}{FastText: Enriching with Subword Information 🔤}
\textbf{The Problem with Word2Vec and GloVe:}

\vspace{1em}

\begin{exampleblock}{Out-of-Vocabulary (OOV) Problem}
\begin{itemize}
    \item Word2Vec/GloVe: One vector per word
    \item New word? \alert{No vector!}
    \item Misspelling? \alert{No vector!}
    \item Rare morphological form? \alert{No vector!}
\end{itemize}
\end{exampleblock}

\vspace{1em}

\textbf{Example:}
\begin{itemize}
    \item Have: "running", "runner"
    \item Need: "runnable" $\rightarrow$ \alert{OOV!}
    \item But these words share morphology! ("run" + suffix)
\end{itemize}

\vspace{1em}

\begin{center}
\Large
\textbf{FastText Solution:} Represent words as \alert{bags of character n-grams}
\end{center}

\small
\textit{Reference: Bojanowski et al. (2017). "Enriching Word Vectors with Subword Information"}
\end{frame}

% ============================================================
% FASTTEXT DETAILS
% ============================================================
\begin{frame}{FastText: Character N-grams 🧩}
\textbf{Key Idea:} Break words into character n-grams

\vspace{1em}

\begin{exampleblock}{Example: "where" (with n=3)}
\textbf{Character trigrams:} <wh, whe, her, ere, re>

\textbf{Plus:} <where> (the whole word)

\textbf{Word vector:}
\begin{align*}
\vec{\text{where}} = \frac{1}{6}\left(\vec{\text{<wh>}} + \vec{\text{whe}} + \vec{\text{her}} + \vec{\text{ere}} + \vec{\text{re>}} + \vec{\text{<where>}}\right)
\end{align*}
\end{exampleblock}

\vspace{1em}

\begin{columns}[T]
\begin{column}{0.48\textwidth}
\textbf{Advantages:}
\begin{itemize}
    \item Handles OOV words!
    \item Captures morphology
    \item Better for rare words
    \item Great for German, Turkish, Finnish
    \item Robust to typos
\end{itemize}
\end{column}

\begin{column}{0.48\textwidth}
\textbf{Example OOV:}

Have trained on: "run", "running"

Need vector for: "runnable"

N-grams: <ru, run, unn, nna, nab, abl, ble, le>

Many overlap with "run" and "running"!

$\rightarrow$ Can generate reasonable embedding
\end{column}
\end{columns}
\end{frame}

% ============================================================
% FASTTEXT CODE
% ============================================================
\begin{frame}[fragile]{FastText in Practice 💻}
\begin{lstlisting}[language=Python]
from gensim.models import FastText
import gensim.downloader as api

# Load pre-trained FastText model
# Available in 157 languages!
fasttext_model = api.load('fasttext-wiki-news-subwords-300')

# Works for known words
vec_cat = fasttext_model['cat']

# Also works for unknown words! (OOV)
vec_unknownword = fasttext_model['unknownword']  # Still get a vector!

# Even works for misspellings (somewhat)
vec_misspelling = fasttext_model['computr']  # Close to "computer"

# Train your own FastText model
sentences = [['the', 'cat', 'sat'], ['the', 'dog', 'ran']]

ft_model = FastText(
    sentences,
    vector_size=100,
    window=5,
    min_count=1,
    min_n=3,      # minimum n-gram length
    max_n=6,      # maximum n-gram length
    word_ngrams=1 # use word + ngrams
)

# Check similar words
similar = ft_model.wv.most_similar('cat', topn=5)

# Morphological awareness
# 'running', 'runner', 'runnable' will be close
\end{lstlisting}
\end{frame}

% ============================================================
% COMPARISON TABLE
% ============================================================
\begin{frame}{Embedding Methods Comparison 📊}
\begin{center}
\small
\begin{tabular}{lccccc}
\toprule
\textbf{Method} & \textbf{Year} & \textbf{Type} & \textbf{OOV} & \textbf{Speed} & \textbf{Best For} \\
\midrule
LSA & 1990 & Count & ✗ & Medium & IR, exploration \\
LDA & 2003 & Probabilistic & ✗ & Slow & Topic modeling \\
Word2Vec & 2013 & Neural (local) & ✗ & Fast & General NLP \\
GloVe & 2014 & Hybrid (global) & ✗ & Fast & Large corpora \\
FastText & 2017 & Neural + ngrams & ✓ & Fast & Morphology-rich \\
\bottomrule
\end{tabular}
\end{center}

\vspace{1em}

\begin{block}{When to Use What?}
\begin{itemize}
    \item \textbf{Word2Vec (Skip-gram):} General purpose, good semantic quality, most popular
    \item \textbf{GloVe:} Large corpus, want deterministic training, good interpretability
    \item \textbf{FastText:} Morphologically rich languages, need OOV handling, many rare words
    \item \textbf{LSA/LDA:} Topic modeling, document similarity, interpretable dimensions
\end{itemize}
\end{block}

\vspace{0.5em}
\textbf{Typical dimensions:} 50-300 (100-300 most common)
\end{frame}

% ============================================================
% EVALUATION
% ============================================================
\begin{frame}{Evaluating Word Embeddings 📏}
\begin{columns}[T]
\begin{column}{0.48\textwidth}
\textbf{Intrinsic Evaluation:}

\vspace{0.5em}
\textbf{1. Word Similarity}
\begin{itemize}
    \item WordSim-353 dataset
    \item SimLex-999 dataset
    \item Spearman correlation with human judgments
\end{itemize}

\vspace{0.5em}
\textbf{2. Word Analogies}
\begin{itemize}
    \item Google Analogy Dataset (19k questions)
    \item Semantic: \textit{Athens:Greece::Baghdad:Iraq}
    \item Syntactic: \textit{good:better::bad:worse}
    \item Accuracy: \% correct
\end{itemize}

\vspace{0.5em}
\textbf{3. Categorization}
\begin{itemize}
    \item Cluster words
    \item Compare to gold categories
\end{itemize}
\end{column}

\begin{column}{0.48\textwidth}
\textbf{Extrinsic Evaluation:}

Use embeddings as features in downstream tasks:

\vspace{0.5em}
\begin{itemize}
    \item \textbf{Sentiment Analysis}
    \item \textbf{Named Entity Recognition}
    \item \textbf{POS Tagging}
    \item \textbf{Machine Translation}
    \item \textbf{Question Answering}
    \item \textbf{Text Classification}
\end{itemize}

\vspace{1em}
\begin{alertblock}{Key Insight}
Intrinsic scores don't always correlate with extrinsic performance!

Best approach: Evaluate on your actual task.
\end{alertblock}
\end{column}
\end{columns}
\end{frame}

% ============================================================
% LIMITATIONS
% ============================================================
\begin{frame}{Limitations of Static Embeddings ⚠️}
\textbf{The fundamental problem: One vector per word type}

\vspace{1em}

\begin{exampleblock}{Example: Polysemy}
\begin{enumerate}
    \item "I deposited money at the \textbf{bank}" (financial institution)
    \item "We sat by the river \textbf{bank}" (riverside)
    \item "The plane will \textbf{bank} left" (tilt)
\end{enumerate}
All get the \alert{same vector}! This conflates different meanings.
\end{exampleblock}

\vspace{1em}

\begin{columns}[T]
\begin{column}{0.48\textwidth}
\textbf{Other Limitations:}
\begin{itemize}
    \item No compositional semantics
    \item "hot dog" $\neq$ "hot" + "dog"
    \item Bias in embeddings
    \item Limited to training corpus
    \item No understanding of negation
    \item Missing multimodal grounding
\end{itemize}
\end{column}

\begin{column}{0.48\textwidth}
\textbf{The Solution:}
\begin{itemize}
    \item \alert{Contextual embeddings!}
    \item Different vector per occurrence
    \item Consider sentence context
    \item ELMo, BERT, GPT, ...
    \item Coming next lecture!
\end{itemize}

\vspace{0.5em}
\begin{center}
% \includegraphics[width=0.8\textwidth]{transformer_preview.png}
\textit{(Image placeholder - transformer architecture preview)}
\end{center}
\end{column}
\end{columns}
\end{frame}

% ============================================================
% BIASES
% ============================================================
\begin{frame}{Bias in Word Embeddings ⚖️}
\textbf{Embeddings learn the biases in their training data}

\vspace{1em}

\begin{exampleblock}{Gender Bias Examples}
\begin{itemize}
    \item $\vec{\text{man}} : \vec{\text{programmer}} :: \vec{\text{woman}} : \vec{\text{homemaker}}$
    \item $\vec{\text{father}} : \vec{\text{doctor}} :: \vec{\text{mother}} : \vec{\text{nurse}}$
    \item "man" more associated with "career", "woman" with "family"
\end{itemize}
\end{exampleblock}

\vspace{0.5em}

\textbf{Other Biases:}
\begin{itemize}
    \item Racial/ethnic stereotypes
    \item Age bias
    \item Religious bias
    \item Socioeconomic bias
\end{itemize}

\vspace{1em}

\textbf{Why This Matters:}
\begin{itemize}
    \item Embeddings used in hiring tools, search, recommendations
    \item Can perpetuate and amplify societal biases
    \item May violate fairness and anti-discrimination principles
\end{itemize}

\vspace{0.5em}
\small
\textit{References: Bolukbasi et al. (2016). "Man is to Computer Programmer as Woman is to Homemaker?"}

\textit{Caliskan et al. (2017). "Semantics derived automatically from language corpora contain human-like biases"}
\end{frame}

% ============================================================
% DEBIASING
% ============================================================
\begin{frame}{Debiasing Approaches 🔧}
\textbf{How can we reduce bias in embeddings?}

\vspace{1em}

\begin{columns}[T]
\begin{column}{0.48\textwidth}
\textbf{1. Post-processing:}
\begin{itemize}
    \item Identify bias subspace
    \item Neutralize: Remove bias from neutral words
    \item Equalize: Ensure equal distances
\end{itemize}

\vspace{0.5em}
\textbf{2. Training-time:}
\begin{itemize}
    \item Counterfactual data augmentation
    \item Adversarial debiasing
    \item Constrained optimization
\end{itemize}

\vspace{0.5em}
\textbf{3. Data curation:}
\begin{itemize}
    \item Balanced corpora
    \item Careful source selection
    \item Remove problematic content
\end{itemize}
\end{column}

\begin{column}{0.48\textwidth}
\textbf{Challenges:}
\begin{itemize}
    \item Hard to define "bias"
    \item May harm performance
    \item Bias is multidimensional
    \item Not a complete solution
    \item Ethical considerations
\end{itemize}

\vspace{1em}
\begin{alertblock}{Important}
Debiasing is an active research area. No perfect solution yet!

Best practice:
\begin{itemize}
    \item Be aware of biases
    \item Document training data
    \item Test for bias
    \item Use multiple evaluation metrics
    \item Consider societal impact
\end{itemize}
\end{alertblock}
\end{column}
\end{columns}
\end{frame}

% ============================================================
% PRACTICAL TIPS
% ============================================================
\begin{frame}{Practical Tips for Using Embeddings 💡}
\begin{enumerate}
    \item \textbf{Start with pre-trained models}
    \begin{itemize}
        \item Word2Vec: Google News (300d, 3B words)
        \item GloVe: Common Crawl (300d, 840B tokens)
        \item FastText: Available in 157 languages
    \end{itemize}

    \item \textbf{Fine-tune if you have domain data}
    \begin{itemize}
        \item Medical: PubMed, clinical notes
        \item Legal: case law, contracts
        \item Social media: tweets, posts
        \item Can significantly improve performance
    \end{itemize}

    \item \textbf{Choose dimensions wisely}
    \begin{itemize}
        \item More dims = more expressiveness, more data needed
        \item 50-100d: small datasets, fast computation
        \item 200-300d: large datasets, better quality
    \end{itemize}

    \item \textbf{Use cosine similarity}
    \begin{itemize}
        \item $\text{similarity}(u,v) = \frac{u \cdot v}{||u|| \cdot ||v||}$
        \item Not Euclidean distance (direction matters more than magnitude)
    \end{itemize}

    \item \textbf{Be aware of biases and limitations}
    \begin{itemize}
        \item Test for bias in your use case
        \item Static embeddings can't handle polysemy
        \item Consider contextual embeddings for better performance
    \end{itemize}
\end{enumerate}
\end{frame}

% ============================================================
% APPLICATIONS
% ============================================================
\begin{frame}{Real-World Applications 🌟}
\begin{columns}[T]
\begin{column}{0.48\textwidth}
\textbf{Search \& Information Retrieval:}
\begin{itemize}
    \item Semantic search
    \item Query expansion
    \item Document ranking
    \item Question answering
\end{itemize}

\vspace{0.5em}
\textbf{Text Classification:}
\begin{itemize}
    \item Sentiment analysis
    \item Spam detection
    \item Topic classification
    \item Intent detection
\end{itemize}

\vspace{0.5em}
\textbf{Sequence Labeling:}
\begin{itemize}
    \item Named entity recognition
    \item POS tagging
    \item Chunking
\end{itemize}
\end{column}

\begin{column}{0.48\textwidth}
\textbf{Recommendation Systems:}
\begin{itemize}
    \item Product recommendations
    \item Content recommendations
    \item User similarity
\end{itemize}

\vspace{0.5em}
\textbf{Machine Translation:}
\begin{itemize}
    \item Word alignment
    \item Transfer learning
    \item Multilingual models
\end{itemize}

\vspace{0.5em}
\textbf{Creative Applications:}
\begin{itemize}
    \item Poetry generation
    \item Analogy discovery
    \item Semantic exploration
    \item Visualization
\end{itemize}
\end{column}
\end{columns}

\vspace{1em}
\textbf{Example:} Google Search uses embeddings to understand query intent and match relevant documents!
\end{frame}

% ============================================================
% DISCUSSION
% ============================================================
\begin{frame}{Discussion Question 💬}
\begin{center}
\Large
\textbf{Why do word analogies work so well?}
\end{center}

\vspace{1em}

\textbf{The famous example:}
\begin{center}
$\vec{\text{king}} - \vec{\text{man}} + \vec{\text{woman}} \approx \vec{\text{queen}}$
\end{center}

\vspace{1em}

\textbf{Think about:}
\begin{itemize}
    \item What does vector subtraction represent linguistically?
    \item Why does this capture semantic relationships?
    \item What are the limitations of this approach?
    \item Does this mean embeddings "understand" gender?
\end{itemize}

\vspace{1em}

\begin{alertblock}{Deeper Question}
Are these embeddings truly capturing \textit{meaning}, or just statistical patterns?

Does the distributional hypothesis have limits?
\end{alertblock}

\small
\textit{Reference: Boleda (2020). "Distributional Semantics and Linguistic Theory"}
\end{frame}

% ============================================================
% SUMMARY
% ============================================================
\begin{frame}{Summary 🎯}
\textbf{What we learned today:}

\begin{enumerate}
    \item \textbf{Word2Vec (2013):} Neural revolution in embeddings
    \begin{itemize}
        \item CBOW: Predict center from context
        \item Skip-gram: Predict context from center
        \item Negative sampling for efficiency
    \end{itemize}

    \item \textbf{GloVe (2014):} Combining count + prediction
    \begin{itemize}
        \item Global co-occurrence statistics
        \item Factorization objective
        \item Ratio of probabilities
    \end{itemize}

    \item \textbf{FastText (2017):} Subword information
    \begin{itemize}
        \item Character n-grams
        \item Handles OOV words
        \item Captures morphology
    \end{itemize}

    \item \textbf{Evaluation:} Intrinsic vs. extrinsic

    \item \textbf{Limitations:} Static, biased, no polysemy

    \item \textbf{Next:} Contextual embeddings solve polysemy!
\end{enumerate}
\end{frame}

% ============================================================
% REFERENCES
% ============================================================
\begin{frame}{Key References 📚}
\small
\textbf{Foundational Papers:}
\begin{itemize}
    \item Mikolov et al. (2013a). "Efficient Estimation of Word Representations in Vector Space"
    \item Mikolov et al. (2013b). "Distributed Representations of Words and Phrases and their Compositionality"
    \item Pennington et al. (2014). "GloVe: Global Vectors for Word Representation"
    \item Bojanowski et al. (2017). "Enriching Word Vectors with Subword Information"
\end{itemize}

\vspace{0.5em}
\textbf{Bias \& Ethics:}
\begin{itemize}
    \item Bolukbasi et al. (2016). "Man is to Computer Programmer as Woman is to Homemaker?"
    \item Caliskan et al. (2017). "Semantics derived automatically from language corpora contain human-like biases"
\end{itemize}

\vspace{0.5em}
\textbf{Theory:}
\begin{itemize}
    \item Boleda, G. (2020). "Distributional Semantics and Linguistic Theory"
    \item Levy \& Goldberg (2014). "Neural Word Embedding as Implicit Matrix Factorization"
\end{itemize}

\vspace{0.5em}
\textbf{Tools:}
\begin{itemize}
    \item Gensim: \url{https://radimrehurek.com/gensim/}
    \item Pre-trained vectors: \url{https://github.com/RaRe-Technologies/gensim-data}
\end{itemize}
\end{frame}

% ============================================================
% QUESTIONS
% ============================================================
\begin{frame}{Questions? 🙋}
\begin{center}
\LARGE
\textbf{Next Lecture:}

\vspace{1em}
Contextual Embeddings: ELMo, Universal Sentence Encoder, BERT

\vspace{2em}
\textit{One word, multiple meanings!}
\end{center}
\end{frame}

\end{document}
